\section{Solver calls, verification and amplification}\label{si:sec:transfer-proof}

This section and the next prove Theorems~\ref{d:thm:algorithm-transfer}
and~\ref{c:thm:transfer}. An algorithm for every Maxwell instance in
the constructed class can be applied to the encoding of a sparse linear system, after
which the decoder recovers the algebraic solution. The resource bounds of
the realization theorems hold on every random draw, while conditioning
and correctness hold on a common good event; the arguments below turn
this into a randomized algorithm with verified output.

\subsection{Rational inputs, capped calls, and verification}
For rational inputs, encode, cap the solver at its known good-instance
bound, and decode. On the common event, $N=\tO(m^{3/2})$, $B_{\rm out}=O(\ell)$,
and both $\kappa_X$ and $\delta^{-1}$ are polynomial in $m$.
Substitution in~\eqref{c:eq:solverbound} gives the second term of
\eqref{c:eq:transfer}; preparation and decoding give the first.
A bad draw is not allowed an unbounded solver call: a timeout or an
invalid output format is rejected.

Verify a candidate using the original sparse residual with outward
short-dyadic arithmetic. This costs $O(m\operatorname{polylog}m)$,
without constructing a common denominator for all rational entries.
For an explicit acceptance rule, put $\nu_M=\|b\|_\infty$ and enclose
$\rho_M=\|Mx-b\|^2/\nu_M^2$ and $b_M=\|b\|^2/\nu_M^2$. Their intermediate
magnitudes are bounded by an explicit polynomial from the input and
fixed output ranges. If that bound is $P_{\rm mag}\ge1$, rounding every one
of the $O(m)$ additions and products on a grid at most
$\eta_0^2/(2^{20}m^2P_{\rm mag}^4)$ gives interval widths at most
$\eta_0^2/64$ for $\rho_M$ and $1/64$ for $b_M$; the elementary product
error bound proves this by induction over the operations. Accept
only when the upper bound for $\rho_M$ is at most $\eta_0^2$ times the
lower bound for $b_M$. Acceptance is sound. The successful decoder
has residual less than $\eta_0\|b\|/4$ by
Remark~\ref{si:rem:continuum-compiler}, so it passes with this margin and
$b_M\ge1$.

Fix a positive rational lower bound for the solver's uniform success
probability, as part of its specification. A good draw and successful
call then have a fixed positive joint probability. Independent capped
trials and the sound acceptance test give inverse-polynomial failure
with $O(\log m)$ repetitions, preserving~\eqref{c:eq:transfer}. The
finite-event quantifier is essential here: correctness does not depend
on which admissible residual field the solver chose.

\subsection{Randomization and verifiable acceptance}

The Maxwell reduction has resource bounds on every random draw, while
its conditioning and correctness guarantees hold on a common good
event. A bad draw may be singular or poorly conditioned, so a runtime
bound for well-conditioned instances cannot control an unrestricted
call on such a draw.

We run each Maxwell solve with a cap equal to a proved upper bound for
a good input of the relevant size and input class. A timeout is rejected.
For a returned field, we decode the saddle coordinates and evaluate
the residual of the original, unrounded real saddle system. This costs
$O(m)$ Real-RAM operations. An accepted result satisfies the desired
projection-relative error regardless of whether its physical draw was
good. Good draws and successful solver calls pass with a fixed positive
probability, after allocating the rounding margins in
Section~\ref{d:app:real-transfer}. Independent capped repetitions amplify
success, and $O(\log m)$ attempts yield any prescribed inverse-polynomial
failure bound. Their cost remains~\eqref{d:eq:sparse-composition-time}.

\section{A transfer from rectangular real inputs}
\label{d:app:real-transfer}

Theorem~\ref{d:thm:main} starts with a short rational square system.
Assumption~\ref{d:ass:slth} permits rectangular matrices and
unrestricted real entries. We bridge these models explicitly.

After removing empty rows, the active row count is at most $m$. Normalize
$F=F_0/(ma_F)$ with $a_F=\max|F_{0,ij}|$ and compute $b_F=F^\top b$ with exact Real-RAM
operations before rounding either the matrix or the source. If $b_F=0$,
then $\Pi_{F_0}b=0$ and $x=0$ is an exact answer. Otherwise let
$\nu_F=\norm{b_F}_\infty$, $g_F=b_F/\nu_F$, and construct the sparse saddle system
\begin{equation}
 K_s\begin{pmatrix}y_{\rm s}\\x_{\rm s}\end{pmatrix}=c_{\rm s},
 \qquad
 K_s=\begin{pmatrix}I&F\\F^\top&0\end{pmatrix},
 \qquad c_{\rm s}=\begin{pmatrix}0\\-g_F\end{pmatrix}.
 \label{d:eq:complexity-saddle}
\end{equation}
Its entries are bounded by one, its dimension is at most $2m$, and it has
at most $3m$ nonzeros. The normal matrix $F^\top F$ is used to analyze
the system, but need not be formed by the algorithm.

Let $P=m^{\lceil c_\kappa\rceil+1}$. The inverse of $K_s$ has norm
at most $4P^2$. A relative saddle residual $\eta_s$ implies a
projection-relative error at most $2P\eta_s$ after rescaling the $x_{\rm s}$
coordinates. Taking $\eta_s=m^{-(6\lceil c_\kappa\rceil+3)}$ therefore suffices for the
squared-error target~\eqref{d:eq:slth-target}.

We round the bounded saddle coefficients and normalized source to
$O_{c_\kappa}(\log m)$ bits, with an explicit perturbation budget that preserves
both inverse bounds and the original saddle residual. An additional
$m$-dimensional identity block ensures that the rounded input's own
nonzero count lies between $m$ and $4m$. This counted padding makes its
condition and accuracy bounds fixed polynomials in its own input
size. The rounded square system is then a valid input to
Theorem~\ref{d:thm:main}. The subsections below supply the
normalization, residual identity, rounding inequalities, and parameter
choices.

Computing $b_F$ before rounding matters when it is small because of
cancellation. Rounding $b$ relative to $\norm{b}_2$ would not in general
control the much smaller projected source. Here $g_F$ is normalized only
after the exact projection-relevant quantity has been formed. The scale
factors retained for final output are Real-RAM words; they do not become
large-bit physical coefficients. This is a Real-RAM reduction for
unrestricted real inputs and a finite-bit reduction after the stated
rounding step.

We prove Theorem~\ref{d:thm:algorithm-transfer}, including the error
criterion and randomized runtime model. The source problem and its
projection-relative target are those in
Assumption~\ref{d:ass:slth}. The same bridge, with
Theorem~\ref{c:thm:main} in place of Theorem~\ref{d:thm:main} and the
moment-message decoder of Section~\ref{c:sec:recovery}, proves the
Real-RAM part of Theorem~\ref{c:thm:transfer}.

\subsection{Input access and exact source normalization}

Let $F_0\in\R^{n_0\times n}$ have full column rank, $m\ge2$
nonzeros, and condition number at most $m^{c_\kappa}$. A sparse nonzero list
identifies all nonempty rows. Delete empty rows and relabel the active
indices in $O(m\log m)$ operations. The remaining row count $n_1$
satisfies $n\le n_1\le m$. Empty rows contribute zero to the column
space and the projected right-hand side, and deleting them does not
change the nonzero singular values. We access $b$ only at active row
indices. This uses the random-access input model of the assumption;
a sequential representation containing additional empty rows would
require its reading cost to be charged separately.

Set
\begin{equation}
 a_F=\max_{i,j}|F_{0,ij}|,\quad F=F_0/(ma_F),\quad b_F=F^\top b.
\end{equation}
Compute $b_F$ with exact Real-RAM operations before any rounding. Let
$P=m^{\lceil c_\kappa\rceil+1}$. Then
$\norm{F}\le1$ and $\sigma_{\min}(F)\ge P^{-1}$.
If $b_F=0$, full column rank gives $\Pi_{F_0}b=0$, so $x=0$ is exact.
Otherwise let $\nu_F=\norm{b_F}_\infty>0$ and $g_F=b_F/\nu_F$.
Construct $K_s,c_{\rm s}$ as in~\eqref{d:eq:complexity-saddle}.
Every entry has absolute value at most one,
$1\le\norm{c_{\rm s}}=\norm{g_F}\le\sqrt n$, the dimension is at most $2m$,
and the nonzero count is at most $3m$.

The exact saddle solution is
$(-Fx_{\rm s*},x_{\rm s*})$, where $F^\top Fx_{\rm s*}=g_F$, so
\begin{equation}
 x_{\rm s*}=F^\dagger b/\nu_F,
 \qquad x_*=(\nu_F/(ma_F))x_{\rm s*}=F_0^\dagger b.
 \label{d:eq:app-saddle-exact}
\end{equation}
The algorithm never forms $F^\top F$. It retains $ma_F,\nu_F$ as Real-RAM
words for final scaling; they are not short-rational physical input
coefficients.

\subsection{Saddle conditioning and the target error}

With $\Pi_{F_0}=F(F^\top F)^{-1}F^\top$, direct multiplication verifies
\begin{equation}
 K_s^{-1}=
 \begin{pmatrix}
 I-\Pi_{F_0}&F(F^\top F)^{-1}\\
 (F^\top F)^{-1}F^\top&-(F^\top F)^{-1}
 \end{pmatrix}.
 \label{d:eq:app-saddle-inverse}
\end{equation}
The four blocks have norm at most $1,P,P,P^2$, respectively. Hence
\begin{equation}
 \norm{K_s^{-1}}\le1+2P+P^2\le4P^2=:Q_s,
 \qquad \norm{K_s}\le2.
 \label{d:eq:app-saddle-bounds}
\end{equation}
For an approximate solution $(y_{\rm s},x_{\rm s})$, write
$\rho_1=y_{\rm s}+Fx_{\rm s}$ and $\rho_2=F^\top y_{\rm s}+g_F$. Eliminating $y_{\rm s}$ gives the exact
identity
\begin{equation}
 F(x_{\rm s}-x_{\rm s*})=\Pi_{F_0}\rho_1-F(F^\top F)^{-1}\rho_2.
 \label{d:eq:app-saddle-residual-identity}
\end{equation}
Therefore $\norm{F(x_{\rm s}-x_{\rm s*})}\le2P\norm{(\rho_1,\rho_2)}$.
Since $g_F=F^\top Fx_{\rm s*}$ and $\norm{F}\le1$,
$\norm{g_F}\le\norm{Fx_{\rm s*}}$. If the relative saddle residual is at
most $\eta_s$, the decoded vector $x=(\nu_F/(ma_F))x_{\rm s}$ satisfies
\begin{equation}
 \frac{\norm{F_0x-\Pi_{F_0}b}}{\norm{\Pi_{F_0}b}}
 =\frac{\norm{F(x_{\rm s}-x_{\rm s*})}}{\norm{Fx_{\rm s*}}}
 \le2P\eta_s.
 \label{d:eq:app-projection-relative}
\end{equation}
No factor $\norm{b}/\norm{\Pi_{F_0}b}$ occurs.

Choose $\eta_s=m^{-(6\lceil c_\kappa\rceil+3)}$. For $m\ge2$,
\begin{equation}
 2P\eta_s=2m^{-5\lceil c_\kappa\rceil-2}\le m^{-5\lceil c_\kappa\rceil}
                 \le n^{-5\lceil c_\kappa\rceil}\le n^{-5c_\kappa}.
\end{equation}
Squaring~\eqref{d:eq:app-projection-relative} proves exactly
\eqref{d:eq:slth-target}. Integer exponents are used to make the
parameters computable with basic arithmetic operations.

\subsection{Rounding to the square theorem's input class}

Choose a dyadic spacing
\begin{equation}
 u_s\le\frac{\eta_s}{1024mQ_s}.
\end{equation}
Round each existing nonzero of $K_s$ to that grid, using identical
roundings for the two occurrences of each $F$ entry, and round $g_F$ on
the same grid. Existing zeros remain zero. Write the resulting matrix
and source as $\widehat K_s,\widehat c_{\rm s}$. A conservative count of
the changed entries gives
\begin{align}
 \norm{\widehat K_s-K_s}
   &\le2mu_s\le\frac{\eta_s}{512Q_s},\\
 \norm{\widehat c_{\rm s}-c_{\rm s}}
   &\le mu_s\le\frac{\eta_s}{1024Q_s}
         \le\frac{\eta_s}{1024}\norm{c_{\rm s}}.
 \label{d:eq:app-saddle-rounding}
\end{align}
By the Neumann series and~\eqref{d:eq:app-saddle-bounds},
$\norm{\widehat K_s^{-1}}\le2Q_s$ and
$\norm{\widehat K_s}\le3$.

Suppose a returned vector $(y_{\rm s},x_{\rm s})$ satisfies
\begin{equation}
 \norm{\widehat K_s(y_{\rm s},x_{\rm s})-\widehat c_{\rm s}}
 \le\frac{\eta_s}{64}\norm{\widehat c_{\rm s}}.
\end{equation}
Then
\begin{equation}
 \norm{(y_{\rm s},x_{\rm s})}
 \le2Q_s(1+\eta_s/64)(1+\eta_s/1024)\norm{c_{\rm s}}
 <3Q_s\norm{c_{\rm s}}.
\end{equation}
Restoring the unrounded coefficients therefore gives
\begin{align}
 \norm{K_s(y_{\rm s},x_{\rm s})-c_{\rm s}}
 &\le\norm{\widehat K_s(y_{\rm s},x_{\rm s})-\widehat c_{\rm s}}
       +\norm{\widehat K_s-K_s}\norm{(y_{\rm s},x_{\rm s})}
       +\norm{\widehat c_{\rm s}-c_{\rm s}}\\
 &\le\eta_s\left(
       \frac{1+1/1024}{64}+\frac{3}{512}+\frac{1}{1024}
       \right)\norm{c_{\rm s}}
 <\frac{\eta_s}{16}\norm{c_{\rm s}}.
 \label{d:eq:app-unrounded-margin}
\end{align}
This margin permits any additional finite-precision decoding budget
to be allocated explicitly before testing the original residual.

Append an $m$ by $m$ identity block with zero source. The padded
matrix has nonzero count $m'$ between $m$ and $4m$, dimension at most
$3m$, and
\begin{equation}
 \kappa_2(\widehat K_s\oplus I_m)
 \le6Q_s=24m^{2\lceil c_\kappa\rceil+2}\le(m')^{2\lceil c_\kappa\rceil+7}.
\end{equation}
Its requested relative accuracy $\eta_s/64$ is at least
$(m')^{-(6\lceil c_\kappa\rceil+9)}$. Thus its condition and accuracy bounds have the
fixed-polynomial form required by the square realization theorem,
expressed in terms of the rounded input's own sparsity. All entries
have $O_{c_\kappa}(\log m)$ bits, and the padding costs $O(m)$.

The dyadic approximations can be computed in the basic Real RAM by
bisection and comparisons on $[-1,1]$, using $O(\log m)$ operations
per entry. No floor operation on unrestricted reals is needed.
Repeated doubling chooses the spacing using the explicit integer
exponent bounds. After this step, the Maxwell encoder operates on
short rational data with its usual bit-operation bounds.

\subsection{Verification, time caps, and amplification}

The decoded square-system answer has a prescribed finite rational
representation. Discard identity-padding coordinates and test the
original residual $\norm{K_s(y_{\rm s},x_{\rm s})-c_{\rm s}}\le\eta_s\norm{c_{\rm s}}$ using exact
Real-RAM operations. Comparing squared norms avoids a square-root
primitive. This test costs $O(m)$. Passing the test implies
\eqref{d:eq:slth-target}; the good-output bound
\eqref{d:eq:app-unrounded-margin} supplies substantial acceptance margin.
Final output uses $n$ multiplications by $\nu_F/(ma_F)$.

For each physical draw, apply a time cap equal to the solver's known
runtime bound for a good Maxwell instance at the relevant parameters. Reject
a timeout or an output that fails the prescribed interface or residual
check. Finite-bit subprograms use a private bounded address region.
Check each proposed index before storing or dereferencing it, and abort
the trial if it exceeds the prescribed polynomial address bound.
Retain the original input and verifier data outside that region; clear
the touched private cells on exit. The tests and cleanup are charged
to the capped trial. On executions within the input class the index
bounds make every guard inactive. On other executions the guards prevent a time cap from
being mistaken for a bound on the bit length of an index.
A good draw has probability at least $1-1/m$, and a successful
solver call on it has a fixed positive probability. Independent
repetition therefore gives a fixed positive acceptance probability
per trial for all sufficiently large inputs. A constant multiple of
$\log m$ capped trials makes the probability of no accepted answer
smaller than any chosen inverse polynomial. An accepted answer is
always valid, because acceptance uses the original real saddle system.
Finitely many small sizes can be handled directly.

The encoder and decoder cost $\Ot((m')^{3/2})=\Ot(m^{3/2})$ on every
draw. On a good draw,
$N=\Ot(m^{3/2})$, $B_{\rm out}=O(\log m)$, and both the condition number and
inverse residual tolerance are polynomial in $m$. Substituting these
quantities into~\eqref{d:eq:maxwell-solver-model} gives
$\Ot(m^{3\theta/2})$ for the capped solve. Verification costs
$O(m)$, and logarithmically many trials do not change the soft-order
bound. We obtain~\eqref{d:eq:sparse-composition-time}, completing the
transfer theorem.

\section{Conditional exponent restrictions}\label{si:sec:hardness}

\begin{corollary}[A conditional exponent restriction]
\label{d:cor:complexity}
If Assumption~\ref{d:ass:slth} holds for some $c_\kappa>0$ and
$\gamma>3/2$, a solver satisfying~\eqref{d:eq:maxwell-solver-model} cannot
exist with $\theta<2\gamma/3$ on the complete constructed Maxwell family.
\end{corollary}

The first term in~\eqref{d:eq:sparse-composition-time} is the encoding and
decoding cost. Thus the implication requires $\gamma>3/2$ even if the
Maxwell solve itself were arbitrarily fast. Under this condition and
$\theta<2\gamma/3$, both polynomial exponents on the right side of
\eqref{d:eq:sparse-composition-time} are strictly below $\gamma$;
logarithmic factors can be absorbed into a smaller positive exponent
gap. This proves the corollary from the transfer theorem.

The hardness exponent must be compatible with upper bounds for the same
input and accuracy model. Exact-arithmetic conjugate gradients on the
normal operator costs $O(mn)\le O(m^2)$, and its conditioning-dependent
convergence bound further restricts the admissible exponent to
$\gamma\le\min\{2,1+c_\kappa\}$; the derivation follows
Corollary~\ref{c:cor:hardness} below. Any stronger applicable upper bound further narrows this
necessary range. Figure~\ref{fig:consequence} of the main text displays the exponent
comparison; values outside an admissible hardness range have no
corresponding lower-bound consequence.



\subsection{The continuum corollary and known upper bounds}
Use exactly the randomized exponent-gap hypothesis of
Assumption~\ref{d:ass:slth}: for fixed $c_\kappa$ and $\gamma$, every
$\widetilde O(m^{\gamma'})$ randomized Real-RAM algorithm,
$\gamma'<\gamma$, fails to solve all instances of
\eqref{c:eq:projectiontarget} with success probability at least $2/3$.
It is a hypothesis, independent of Theorem~\ref{c:thm:main}. Hardness
statements in the unit-cost Real-RAM model depend on the set of unit-cost
operations; here it consists of $+,-,\times,\div$ and comparisons, and
excludes the floor function.

\begin{corollary}
\label{c:cor:hardness}
If that matching hypothesis holds for an admissible exponent
$\gamma>3/2$, a solver satisfying~\eqref{c:eq:solverbound} cannot have
$\theta<2\gamma/3$ on the complete constructed continuum family.
\end{corollary}
\begin{proof}
Under both strict inequalities, $3/2<\gamma$ and
$3\theta/2<\gamma$. The two exponents in~\eqref{c:eq:transfer} are
strictly smaller than $\gamma$; logarithmic factors can be absorbed
in a smaller positive exponent gap. The amplified algorithm has
success above $2/3$, contradicting the stated hypothesis.
\end{proof}

The admissibility qualification cannot be omitted. In the exact-arithmetic
Real-RAM model of Assumption~\ref{d:ass:slth}, conjugate gradients applied to the positive
 definite normal operator $F_0^TF_0$ solves the projected problem
exactly in at most $n$ iterations. Each normal-operator product is
two sparse products and costs $O(m)$; the remaining vector work is
$O(n)\le O(m)$. There is no need to form the normal matrix.
For completeness the recurrence uses only $+,-,\times,\div$:
\[
 a^{\rm cg}_j=\frac{\rho_j^T\rho_j}{p_j^T(F_0^TF_0)p_j},\quad
 x_{j+1}=x_j+a^{\rm cg}_jp_j,\quad
 \rho_{j+1}=\rho_j-a^{\rm cg}_j(F_0^TF_0)p_j,
\]
\[
 b^{\rm cg}_j=\frac{\rho_{j+1}^T\rho_{j+1}}{\rho_j^T\rho_j},\qquad
 p_{j+1}=\rho_{j+1}+b^{\rm cg}_jp_j,
 \quad \rho_0=p_0=F_0^Tb.
\]
Stop when $\rho_j=0$. The usual induction, substituting these identities,
makes the nonzero directions mutually $F_0^TF_0$-orthogonal and
makes the residual orthogonal to their span. Thus each nonterminal
step adds an independent direction, so after at most $n$ steps the
residual is zero. Positivity makes all used denominators nonzero.
This gives $O(mn)\le O(m^2)$ exact Real-RAM operations, including
zero-projection sources. The same exact-arithmetic upper bound is
explicitly recorded in~\citep[p.~3:2]{andoni2019}.

The conditioning assumption gives a further restriction when $c_\kappa<1$.
Let $x_*=F_0^\dagger b$.
The same orthogonality induction shows that the $t$th CG iterate
minimizes the energy error over its Krylov space. Therefore, for
every polynomial $p$ of degree at most $t$ with $p(0)=1$,
\[
 \frac{\|F_0(x_t-x_*)\|}{\|\Pi_{F_0}b\|}
 \le \max_{\lambda\in\spec F_0^TF_0}|p(\lambda)|.
\]
For completeness, the error candidates are $p(F_0^TF_0)x_*$: writing
$p(\lambda)=1-\lambda p_1(\lambda)$ makes the corresponding iterate
$p_1(F_0^TF_0)F_0^Tb$ belong to the Krylov space. Orthogonality to that
space proves the energy-minimization statement by Pythagoras.
On the interval $[\lambda_-,\lambda_+]$ containing the spectrum, use
\[
 p(\lambda)=
 \frac{T_t((\lambda_-+\lambda_+-2\lambda)/(\lambda_+-\lambda_-))}{T_t((\lambda_-+\lambda_+)/(\lambda_+-\lambda_-))}.
\]
The recurrence for Chebyshev polynomials gives $|T_t|\le1$ on
$[-1,1]$ and $T_t(s)=(w^t+w^{-t})/2$ for
$w=s+\sqrt{s^2-1}>1$. Hence, with
$\kappa=\kappa_2(F_0)$,
\[
 \frac{\|F_0(x_t-x_*)\|}{\|\Pi_{F_0}b\|}
 \le2\left(\frac{\kappa-1}{\kappa+1}\right)^t
 \le2e^{-2t/\kappa}.
\]
The case $\kappa=1$ terminates in one step. Thus
$O_{c_\kappa}(m^{c_\kappa}\log(2+m))$ iterations suffice for the target
$n^{-5c_\kappa}$, and finite termination also caps the count at $n$.
No square root is used in the CG algorithm; the displayed polynomials
only analyze it. For an arbitrary fixed real exponent $c_\kappa$,
one may choose any fixed rational $c_\kappa'>c_\kappa$ in a counted iteration
cap. The integer ceiling of $m^{c_\kappa'}$ is computable by integer
bisection on a fixed-power inequality, so this uses only the stated
Real-RAM operations. Taking $c_\kappa'$ arbitrarily close to $c_\kappa$ is enough
for the exponent-gap comparison.

Accordingly, the matching hypothesis is contradicted whenever
$\gamma>\min\{2,1+c_\kappa\}$. Its necessary upper-bound-compatible range
in this corollary is
\begin{equation}
 3/2<\gamma\le\min\{2,1+c_\kappa\},\qquad \theta\ge2\gamma/3
 \quad\hbox{conditionally on the matching hypothesis}.          \label{c:eq:gamma-range}
\end{equation}
For $c_\kappa\le1/2$ that range is empty; for $c_\kappa\ge1$ its displayed
upper endpoint is two. Admissibility also requires consistency
with any stronger algorithmic upper bound in the same model.
No hardness assertion for any member of this range is made here.
The finite-bit algorithms for sparse systems with short encoded
inputs~\cite{pengvempala2021,nie2022}, including the recent
$\widetilde O(mn\log(\kappa/\eta))$ result~\cite{kelner2026}, are
consistent with this restriction; it is not substituted for the unrestricted-real,
projection-relative arithmetic model above. The continuum construction
and its quantitative realization are independent of accepting any
sparse-system hardness hypothesis.

\subsection{Interpretation of the consequence}

The corollary concerns a uniform worst-case solver over the constructed
family. It identifies the computational consequences of supporting every
instance in that family with the stated dependence on conditioning and
accuracy. An algorithm whose cost is $N$ times a polynomial in
$\kappa_2(A)$ does not satisfy~\eqref{d:eq:maxwell-solver-model}: on these
instances the condition number is allowed to grow polynomially with the
input size. The distinction between polynomially bounded instances and
polylogarithmic algorithmic dependence on their condition number must be
retained when stating the result.

The theorem leaves room for faster algorithms on additional restricted
families, for instance when the admissible geometry or the conditioning
assumption is stronger. It also does not measure the time required by an
analog device to prepare a field, settle, or expose the prescribed output
precision. Those questions require a physical resource model in addition
to the linear-system realization proved here.
